\documentclass[10pt,a4paper]{amsart}
\usepackage[margin=19mm]{geometry}
\usepackage{amsmath,amssymb,mathtools}
\usepackage[scr=rsfs,cal=euler]{mathalfa}
\usepackage{xfrac}
\usepackage[unicode,hidelinks,bookmarks=false]{hyperref}
\newcommand{\ie}{i.e.\ }
\newcommand{\cf}{cf.\ }
\newcommand{\resp}{respectively\ }
\newcommand{\Subsection}{\subsection}
\providecommand{\nobreakdash}{\nobreak\hbox{-}}
\setlength{\emergencystretch}{2em}
\multlinegap0pt
\usepackage{fullpage,mathrsfs,mathtools,enumitem,indentfirst}
\usepackage{bbm}
\usepackage{tikz}
\usetikzlibrary{arrows.meta}
\usepackage{braket}
\usepackage{algorithm2e}
\newtheorem{theorem}{Theorem}
\newtheorem{lemma}{Lemma}
\newtheorem{proposition}{Proposition}
\newtheorem{definition}{Definition}
\newtheorem{corollary}{Corollary}
\newtheorem{observation}{Observation}
\newtheorem{example}{Example}
\newtheorem{claim}{Claim}
\newtheorem{fact}{Fact}
\newtheorem{question}{Question}
\newtheorem{remark}[theorem]{Remark}
\newtheorem*{exercise}{Exercise}
\newcommand{\indpt}{\mathrel{\text{\scalebox{1.07}{$\perp\mkern-10mu\perp$}}}}
\newcommand{\diff}{\partial}
\newcommand{\prob}{\mathbb{P}}
\newcommand{\tr}{\text{tr}}
\newcommand{\E}{\mathbb{E}}
\newcommand{\rank}{\mathrm{rank}}
\newcommand{\srank}{\mathrm{srank}}
\newcommand{\CI}{\mathrm{CI}}
\newcommand{\h}{\mathrm{h}}
\newcommand{\TV}{\mathrm{TV}}
\newcommand{\JS}{\mathrm{JS}}
\newcommand{\KL}{\mathrm{KL}}
\newcommand{\ID}{\mathrm{ID}}
\newcommand{\mic}{\textcolor{magenta}}
\newcommand{\nid}{n^*_\mathrm{id}}
\newcommand{\nnonid}{n^*_{\mathrm{non-id}}}
\newcommand{\nseq}{n^*_{\mathrm{seq}}}
\DeclareMathOperator*{\argmax}{arg\,max}
\DeclareMathOperator*{\argmin}{arg\,min}
\newcommand{\skipitems}[1]{%
  \addtocounter{\@enumctr}{#1}%
}
\begin{document}
\title[Counterexample to a Proposed Capacity Characterization of th]{Counterexample to a Proposed Capacity Characterization of the Relay Channel}
\author{Chun Hei Michael Shiu}
\date{}
\maketitle
\noindent\emph{Source and licence.} Counterexample to a Proposed Capacity Characterization of the Relay Channel. Version arXiv:2609.18727v1 (CC BY 4.0). This material is distributed under the Creative Commons Attribution 4.0 International licence, http://creativecommons.org/licenses/by/4.0/, as recorded in the arXiv OAI licence metadata for this exact version and shown on the abstract page https://arxiv.org/abs/2609.18727v1. Full rights notice: licenses/arxiv-2609.18727-CC-BY-4.0.txt.\par\smallskip
\noindent\textit{Editorial scope.} Counted unit: the original \texttt{lemma} environment \texttt{-} at source lines 284--286 together with its complete proof at lines 288--295; the delivered document keeps source lines 191--296 so that every referenced label and every citation resolves inside it. Native editable source is the paper's own concatenated TeX; the AMS wrapper is editorial formatting only. No publisher class, upstream script or unrelated artwork is redistributed.
\par\smallskip\noindent\textit{Reference note.} This article ships no compiled bibliography (biblatex); the entries delivered here are composed from the author's own BibTeX (.bib) fields, with no field invented and no entry dropped.
\begin{theorem}[Theorem 4 in \cite{pon26}] \label{thm: main thm}
    For a relay channel $(\mathcal X_s,\mathcal X_r, \mathcal Y_r, \mathcal Y_d, \{p(y_r, y_d | x_s, x_r)\})$, the capacity of the relay channel is 
    \begin{align*}
        C_{\mathrm{prop}} = \max\{R_{\mathrm{U-CF}}, R_{\mathrm{C-CF}}, R_{\mathrm{DF}}\},
    \end{align*}
    where the uncoordinated-compress-forward rate $R_{\mathrm{U-CF}}$ is defined as
    \begin{align*}
        R_{\mathrm{U-CF}} := \max_{(p(x_s, x_r), p(y_r |\hat{y}_r, x_r), p(\hat{y}_r|x_r)) \in \mathcal P_{\mathrm{ind-feasible}}} I_{\mathrm{cf}}(X_s; \hat{Y}_r, Y_d | X_r);
    \end{align*}
    the coordinated-compress-forward rate $R_{\mathrm{C-CF}}$ is defined as 
    \begin{align*}
        R_{\mathrm{C-CF}} := \max_{(p(x_s, x_r), p(y_r |\hat{y}_r, x_r), p(\hat{y}_r|x_r)) \in \mathcal P_{\mathrm{cor-feasible}}} I_{\mathrm{for}}(X_s; \hat{Y}_r, Y_d | X_r);
    \end{align*}
    and the decode-forward rate $R_{\mathrm{DF}}$ is defined as 
    \begin{align*}
        R_{\mathrm{DF}} := \max_{p(x_s, x_r)} \min\{I(X_s; Y_r | X_r), I(X_s, X_r; Y_d)\}
    \end{align*}
    with joint distribution $p(x_s, x_r, y_r, y_d) = p(x_s, x_r) \cdot p(y_r, y_d | x_s, x_r)$.
\end{theorem}


\section{Counterexample and Proofs} \label{sec: counterexample and proofs}
In this section, we present a relay channel whose capacity is strictly larger than the proposed capacity in Theorem \ref{thm: main thm}, thereby serving as a counterexample to Theorem 4 in \cite{pon26}.% Again, we emphasize that the counterexample is constructed by ChatGPT-6 Astra.


\subsection{Counterexample Relay Channel} \label{sec: counterexample channel}
Consider the following relay channel $\mathcal R$. Let the (sender) source input be $X_s = (A,B) \in \{0,1\}^2$ be a binary pair, and also let the relay input $X_r \in \{0,1\}$ be binary. The relay channel outputs
\begin{align*}
    Y_r = A \oplus N_1, \qquad Y_d = (B, Z) = (B, X_r \oplus N_2),
\end{align*}
where $N_1, N_2 \sim \mathrm{Bern}(1/4)$ are mutually independent noises, independent of the current channel inputs, with fresh independent copies across channel uses. It is clear that $\mathcal X_s = \{0,1\}^2 = \mathcal Y_d$, $\mathcal X_r = \{0,1\} = \mathcal Y_r$. The relay channel $\mathcal R$ is depicted in Figure \ref{fig:counterexample-relay}.

\begin{figure}[ht]
    \centering
    \begin{tikzpicture}[
        >=Stealth,
        line width=0.7pt,
        block/.style={
            draw,
            minimum height=0.85cm,
            align=center,
            inner sep=6pt
        },
        lab/.style={
            midway,
            above=4pt,
            inner sep=1pt
        }
    ]
        % Main blocks
        \node[block, minimum width=1.8cm]
            (encoder) at (-5,0) {Encoder};

        \node[block, minimum width=3.8cm]
            (channel) at (0,0)
            {$p(y_r,(b,z)\mid(a,b),x_r)$};

        \node[block, minimum width=1.8cm]
            (decoder) at (5,0) {Decoder};

        \node[block, minimum width=2.8cm]
            (relay) at (0,2.8) {Relay encoder};

        % Horizontal links
        \draw[->] (-7,0)
            -- node[lab] {$M$} (encoder.west);

        \draw[->] (encoder.east)
            -- node[lab] {$(A,B)$} (channel.west);

        \draw[->] (channel.east)
            -- node[lab] {$(B,Z)$} (decoder.west);

        \draw[->] (decoder.east)
            -- node[lab] {$\hat M$} (7,0);

        % Vertical relay links
        \draw[->] ([xshift=-0.9cm]channel.north)
            -- node[midway,left=5pt] {$Y_r$}
            ([xshift=-0.9cm]relay.south);

        \draw[->] ([xshift=0.9cm]relay.south)
            -- node[midway,right=5pt] {$X_r$}
            ([xshift=0.9cm]channel.north);
    \end{tikzpicture}
    \caption{Counterexample relay channel $\mathcal R$.}
    \label{fig:counterexample-relay}
\end{figure}

\subsection{Lower Bound for the Capacity of $\mathcal R$} \label{sec: lower bound}

In this section, we describe an achievability scheme for the above relay channel, hence illustrating a lower bound for its capacity. 


\begin{lemma}
    Let $\mathcal R$ be the relay channel defined in Section \ref{sec: counterexample channel}. Then $C(\mathcal R) \geq 2 - h_2(1/4)$.
\end{lemma}

\begin{proof}
The component $B$ supports one bit per channel use over the noiseless source-destination link. Independently, the source transmits a sub-message through $A \to Y_r$ at any rate below $1 - h_2(1/4)$. The relay decodes this sub-message from one coding block and forwards it through $X_r \to Z$ in the following block. By pipelining over sufficiently many blocks, the initialization and termination overhead vanishes. Thus, the total rate can be arbitrarily close to $2 - h_2(1/4)$.
% We may treat the two information bits $(A, B)$ separately. Note that the channel $\mathcal R$ transmit the information bit $B$ noiselessly to the receiver (up to a permutation in the coordinate). Therefore, we can always use exactly 1 bit to transmit $B$. On the other hand, the channel $\mathcal R$ can be viewed as sending $A$ to $Y_r$ through a binary symmetric channel $\mathrm{BSC}(1/4)$, then decode-forward and send the relay source $X_r$ to $Z$ through another binary symmetric channel $\mathrm{BSC}(1/4)$. Therefore, the rate for transmitting $A$ is
% \begin{align*}
%     \min\{C_{A\to Y_r}, C_{X_r \to Z}\} = \min\{1- h_2(1/4), 1-h_2(1/4)\} = 1-h_2(1/4)
% \end{align*}
% as the channel capacity of $\mathrm{BSC}(\delta)$ is $1 - h_2(\delta)$. As they are conducted in parallel, the overall rate is $1 + 1 - h_2(1/4) = 2 - h_2(1/4)$.
\end{proof}

\begin{thebibliography}{99}
\bibitem[]{pon26} Jonathan Ponniah. The Capacity of the Relay Channel. (2026)
\end{thebibliography}
\end{document}
